% TEMPLATE-MILC-v3.tex - plantilla maestra UNICA de las guias MILC v3.
%
% La usan las cuatro familias de guias, cada una con su driver:
%   · Grados 6-11          -> scripts/build-guias-g11.py
%   · Bebras               -> scripts/build-guias-bebras.py
%   · Semillero            -> scripts/build-guias-semillero.py
%   · Territorio interior  -> scripts/build-guias-territorio-interior.py
%
% Antes habia tres copias casi identicas (~400 lineas iguales, 6 distintas):
% cada mejora habia que aplicarla tres veces, y en la practica no se hacia
% (el motor de recursos vivia solo en dos de ellas). Lo que varia entre
% familias esta parametrizado en 6 placeholders que cada driver llena
% (se nombran aqui SIN su sintaxis real: el builder reemplaza por texto plano,
% tambien dentro de los comentarios, y un valor multilinea rompe el preambulo):
%
%   HEADER_IZQ        encabezado izquierdo de pagina
%   HEADER_DER        encabezado derecho de pagina
%   PORTADA_ETIQUETA  rotulo sobre el numero grande (GRADO/MOMENTO/MODULO)
%   PORTADA_SERIE     linea de serie de la portada
%   PORTADA_ID        identificador de la guia en la portada
%   PORTADA_CIRCULO   el circulito de grado (°), o vacio si no aplica
%   PDF_TITULO        titulo en texto plano (sin \textbf ni comillas TeX) para
%                     los metadatos del PDF (/Title); lo produce lib_xelatex.texto_pdf
%
% Composicion (auditoria 2026-09-03): las bandas de estacion piden espacio
% (\Needspace*) y prohiben el salto justo despues (\nopagebreak), asi no quedan
% huerfanas al pie; las cajas largas (softbox, drawbox, infoband, puentebox) son
% `breakable`, asi una caja de media pagina ya no arrastra todo a la siguiente
% dejando la anterior al 40 %. Todo texto sobre mostaza o verde va en milcNegro
% (blanco sobre #E5B400 da 1,93:1 y sobre #58A923 2,95:1; ninguno pasa WCAG AA).
% El resto de claves las documenta content/guias/_SCHEMA.md. El script valida
% que no queden placeholders sin reemplazar antes de escribir el archivo final.

% Etiquetado PDF (latex-lab): /Lang, estructura y texto alternativo de las
% figuras, para lectores de pantalla. Debe ir ANTES de \documentclass.
\DocumentMetadata{lang=es-CO,tagging=on}
\documentclass[11pt,letterpaper]{article}
% headheight=24pt: la cabecera derecha lleva dos lineas (identidad de la guia
% y estacion actual). headsep se acorta en la misma medida para que el bloque
% de texto quede exactamente donde estaba con headheight=12pt.
\usepackage[margin=2.0cm,headheight=24pt,headsep=13pt]{geometry}
\usepackage[table]{xcolor}
\usepackage{fontspec}
\usepackage[spanish]{babel}
\usepackage{tikz}
% Librerías TikZ para los diagramas de las guías (bloque `recursos:`):
%  · babel  → babel en español vuelve activos caracteres como «>», y TikZ los usa
%             en sus opciones (>=stealth, ->). Sin esta librería, cualquier
%             diagrama con flechas revienta con «\language@active@arg> has an
%             extra }». Debe ir ANTES de que se dibuje nada.
%  · positioning        → sintaxis «below left=2pt and 3pt».
%  · decorations.pathreplacing → llaves ({brace}) para señalar rangos.
\usetikzlibrary{babel,positioning,decorations.pathreplacing}
\usepackage{graphicx}
% Circuitos: el trabajo de electrónica se hace hoy con micro:bit y sus sensores
% integrados (MakeCode, sin cableado), así que no se necesitan esquemáticos.
% Si algún día entran montajes con protoboard, basta descomentar esta línea:
% los diagramas .tex/.tikz declarados en `recursos:` ya se incrustan solos.
% \usepackage{circuitikz}
\usepackage[most]{tcolorbox}
\usepackage{tabularx}
\usepackage{array}
\usepackage{fancyhdr}
\usepackage{titlesec}
\usepackage[document]{ragged2e}  % bandera a la derecha en todo el cuerpo
\usepackage{enumitem}
\usepackage{microtype}
\usepackage{etoolbox}
\usepackage{needspace}
% hyperref al final del preambulo: metadatos del PDF (titulo, autor, idioma,
% familia), marcadores por estacion (\pdfbookmark) y enlaces sin recuadro.
\usepackage[hidelinks,
  pdftitle={Los grupos de chat: lo que se reenvía sin cara y sin fecha de cierre},
  pdfauthor={PhD. Álvaro Cárdenas Orozco},
  pdfsubject={Territorio interior · Educación socioemocional · Grado 7° · Momento 8 · MILC v3},
  pdflang={es-CO},
  bookmarksopen=true,bookmarksnumbered=false]{hyperref}

% Helvetica Neue no trae la flecha U+2192, asi que cada «→» escrito literalmente
% en el YAML se compilaba a NADA, en silencio (462 flechas en 61 guias: rutas del
% tipo «paleta → tipografia → lockup» salian rotas). Mapeamos el caracter a su
% version matematica, que si tiene glifo. Asi el autor puede seguir escribiendo
% «→» comodamente en el YAML.
\usepackage{newunicodechar}
\newunicodechar{→}{\ensuremath{\rightarrow}}
% Mismo problema con los simbolos de marca y vinetas: el «✓» de «una palomita ✓»
% salia como cuadrito vacio en el PDF de todas las guias que lo usaban.
\usepackage{pifont}
\newunicodechar{✓}{\ding{51}}
\newunicodechar{✗}{\ding{55}}
\newunicodechar{✔}{\ding{52}}
\newunicodechar{•}{\textbullet}
\newunicodechar{≥}{\ensuremath{\geq}}
\newunicodechar{≤}{\ensuremath{\leq}}

\setmainfont{Helvetica Neue}
\setsansfont{Helvetica Neue}
\setlength{\parindent}{0pt}
\setlength{\parskip}{6pt}
% Sin particion de palabras: en 6.º-7.º una palabra cortada por guion al final
% de linea («Comuni-cacion») frena la lectura mucho mas que una linea corta.
\hyphenpenalty=10000
\exhyphenpenalty=10000
\pretolerance=2000
\tolerance=3000
\setlength{\emergencystretch}{3em}
\linespread{1.10}
\renewcommand{\arraystretch}{1.25}
\setlist{leftmargin=5mm,itemsep=1mm,topsep=1mm}
\newcolumntype{Y}{>{\RaggedRight\arraybackslash}X}

\definecolor{milcMagenta}{HTML}{E6007E}
\definecolor{milcVino}{HTML}{5A0038}
\definecolor{milcTurquesa}{HTML}{0093A5}
\definecolor{milcVerde}{HTML}{58A923}
\definecolor{milcMostaza}{HTML}{E5B400}
\definecolor{milcOcre}{HTML}{B86600}
\definecolor{milcNegro}{HTML}{241816}
\definecolor{milcGris}{HTML}{F7F7F4}
\definecolor{milcLinea}{HTML}{E8E2DD}
\definecolor{milcGradePrimary}{HTML}{0D9488}
\definecolor{milcGradeDark}{HTML}{115E59}
\definecolor{milcGradeSoft}{HTML}{CCFBF1}
\definecolor{milcGradeAccent}{HTML}{CFE7FF}
\definecolor{milcGradeText}{HTML}{FFFFFF}
\color{milcNegro}

\pagestyle{fancy}
\fancyhf{}
% Cabecera de dos lineas: identidad de la guia arriba y, a la derecha y en
% gris, la estacion en curso (\markright desde \milcestacion). Antes de la
% primera estacion la segunda linea va vacia; el \strut mantiene la altura
% para que la linea de arriba no baile entre paginas.
\lhead{\small Territorio interior · Educación socioemocional\\[-1pt]{\footnotesize\strut}}
\rhead{\small Grado 7° · Momento 8 · MILC v3\\[-1pt]{\footnotesize\strut\color{milcNegro!65}\rightmark}}
\cfoot{\small\thepage}
\renewcommand{\headrulewidth}{0pt}

% Sobre mostaza (#E5B400) y verde (#58A923) el texto va en milcNegro; sobre el
% resto de la paleta, en blanco. Se usa dentro de fonttitle/fontupper, donde un
% \color posterior manda sobre coltitle.
\newcommand{\milctintasobre}[1]{%
  \ifboolexpr{test{\ifstrequal{#1}{milcMostaza}} or test{\ifstrequal{#1}{milcVerde}}}%
    {\color{milcNegro}}{\color{white}}}

\newtcolorbox{titlebox}[1]{
  enhanced,colback=#1,colframe=#1,arc=7mm,boxrule=0pt,
  left=7mm,right=7mm,top=3mm,bottom=3mm,
  before skip=8pt,after skip=8pt,
  fontupper=\sffamily\bfseries\Large\color{white}
}

\newtcolorbox{softbox}[3]{
  enhanced,breakable,pad at break=3mm,
  colback=#2,colframe=#1,borderline west={4pt}{0pt}{#1},
  arc=2mm,boxrule=.4pt,left=4mm,right=4mm,top=3mm,bottom=3mm,
  before skip=6pt,after skip=8pt,title={#3},coltitle=white,
  colbacktitle=#1,fonttitle=\sffamily\bfseries\milctintasobre{#1},fontupper=\RaggedRight,
  attach boxed title to top left={xshift=3mm,yshift=-2mm},
  boxed title style={arc=2mm, boxrule=0pt}
}

\newtcolorbox{drawbox}[2]{
  enhanced,breakable,pad at break=4mm,
  colback=white,colframe=#1,arc=4mm,boxrule=1pt,
  left=5mm,right=5mm,top=4mm,bottom=4mm,before skip=8pt,after skip=8pt,
  title={#2},coltitle=white,colbacktitle=#1,fonttitle=\sffamily\bfseries\milctintasobre{#1},
  fontupper=\RaggedRight,
  attach boxed title to top left={xshift=4mm,yshift=-2mm},
  boxed title style={arc=3mm, boxrule=0pt}
}

\newtcolorbox{infoband}[1]{
  enhanced,breakable,pad at break=2mm,colback=#1!8,colframe=#1!50,arc=2mm,boxrule=.5pt,
  left=4mm,right=4mm,top=2mm,bottom=2mm,before skip=4pt,after skip=4pt,
  fontupper=\small\RaggedRight
}

% Puente narrativo entre secciones (contrato editorial Conéctate).
% Da continuidad: "Acabas de X. Ahora vas a Y porque Z."
% Visualmente: banda izquierda en color del grado, texto en itálica pequeña.
\newtcolorbox{puentebox}{
  enhanced,breakable,pad at break=2mm,colback=milcGradeSoft,colframe=milcGradeDark,
  borderline west={3pt}{0pt}{milcGradeDark},
  arc=0mm,boxrule=0pt,left=5mm,right=4mm,top=2.5mm,bottom=2.5mm,
  before skip=4pt,after skip=8pt,
  fontupper=\itshape\small\color{milcNegro}\RaggedRight
}
% La flecha va en modo matematico a proposito: Helvetica Neue NO tiene el glifo
% U+2192, asi que \textrightarrow se compilaba a NADA («Missing character: There
% is no → in font Helvetica Neue») y el puente salia sin flecha. En modo
% matematico la flecha viene de la fuente matematica y si se dibuja.
\newcommand{\puente}[1]{%
  \begin{puentebox}%
    \textcolor{milcGradeDark}{$\rightarrow$}\hspace{2mm}#1%
  \end{puentebox}%
}

\newcommand{\linead}[1]{\noindent\color{milcLinea}\rule{#1}{0.5pt}\color{milcNegro}}
\newcommand{\respuesta}[1]{\par\vspace{2mm}\foreach \n in {1,...,#1}{\noindent\color{milcLinea}\rule{\linewidth}{0.5pt}\par\vspace{5mm}}\color{milcNegro}}
\newcommand{\checkbox}{\textcolor{milcGradeDark}{\fbox{\phantom{x}}}\hspace{5pt}}

% ─── Listas aireadas para las guías socioemocionales ────────────────────
% Componer las verificaciones como «\checkbox texto\\» deja el texto que salta
% de línea debajo del cuadrito y apelmaza el bloque. Estos entornos alinean la
% continuación y airean los ítems. Los usa Territorio Interior; cualquier
% familia puede hacerlo.
\newlist{checklista}{itemize}{1}
\setlist[checklista]{
  label=\textcolor{milcGradeDark}{\fbox{\phantom{x}}},
  leftmargin=9mm, labelsep=3.2mm, itemsep=2.4mm,
  topsep=2.5mm, parsep=0pt, partopsep=0pt, before=\RaggedRight
}
\newlist{pasoslista}{enumerate}{1}
\setlist[pasoslista]{
  label=\textcolor{milcGradeDark}{\sffamily\bfseries\arabic*.},
  leftmargin=8mm, labelsep=3mm, itemsep=2.8mm,
  topsep=1mm, parsep=0pt, partopsep=0pt, before=\RaggedRight
}

\newcommand{\phasecard}[4]{
\begin{tcolorbox}[enhanced,colback=#3,colframe=#2,arc=3mm,boxrule=.5pt,height=3.05cm,valign=center,halign=center,left=1.5mm,right=1.5mm,top=2mm,bottom=2mm]
{\sffamily\bfseries\fontsize{22}{24}\selectfont\textcolor{#2}{#1}}\par
{\sffamily\bfseries\small #4}
\end{tcolorbox}
}

% ── Piezas del rediseno 2026 (legibilidad 6.º-11.º) ───────────────────
% Banda de estacion: reemplaza el titlebox macizo. Titulo a la izquierda,
% dato practico (tiempo/modalidad) a la derecha, en una sola linea.
% Nunca huerfana: pide 9 lineas libres (o salta de pagina rellenando con
% \vfill) y prohibe el salto justo despues de la banda. Nueve y no seis:
% la caja partible que sigue necesita ~80pt (titulo adosado, relleno y dos
% lineas) para arrancar en la misma pagina; con menos, tcolorbox la manda
% entera a la siguiente y la banda se queda sola. El penalty 10000 va
% en `after`, ANTES del espacio posterior: un `after skip` (aunque sea 0pt)
% es un \vskip tras la caja, y ese glue es un punto de corte legal que un
% \nopagebreak escrito despues de \end{tcolorbox} ya no cierra. Cada banda
% deja marcador PDF y actualiza la cabecera derecha.
\newcounter{milcestacion}
\newcommand{\milcestacion}[2]{%
\Needspace*{9\baselineskip}%
\stepcounter{milcestacion}%
\pdfbookmark[1]{#2}{milcestacion\themilcestacion}%
\markright{#2}%
\begin{tcolorbox}[enhanced,colback=#1,colframe=#1,arc=2mm,boxrule=0pt,
  left=5mm,right=5mm,top=2.6mm,bottom=2.6mm,before skip=11pt,
  after={\par\nopagebreak\vskip7pt}]
{\sffamily\bfseries\fontsize{15}{18}\selectfont\milctintasobre{#1}#2}
\end{tcolorbox}}

% Tarjeta de paso numerada: sustituye las tablas «Pilar / Aplicacion», que
% eran formato de consulta docente, no de estudiante. El numero va en un
% circulo sobre el borde para que la secuencia se lea de un vistazo.
\newcommand{\milcpaso}[3]{%
\Needspace{4\baselineskip}%
\begin{tcolorbox}[enhanced,colback=white,colframe=milcLinea,arc=1.5mm,boxrule=.9pt,
  left=14mm,right=4mm,top=3mm,bottom=3mm,before skip=4pt,after skip=4pt,
  fontupper=\RaggedRight,
  overlay={\node[anchor=north west,circle,fill=#2,text=white,inner sep=0pt,
    minimum size=7.4mm,font=\sffamily\bfseries\small]
    at ([xshift=3.6mm,yshift=-3.2mm]frame.north west) {#1};}]
#3
\end{tcolorbox}}

% Casilla marcable de tamano util con lapiz (la anterior era \fbox{\phantom{x}},
% de alto variable segun el texto que la seguia).
\newcommand{\milcmarca}{\framebox[3.8mm]{\rule{0pt}{3.4mm}}}

% Fila marcable de lista suelta (checks de la estacion 1).
\newcommand{\milccheckrow}[1]{%
\par\vspace{1.8mm}\noindent
\makebox[6.4mm][l]{\raisebox{-.55mm}{\textcolor{milcNegro}{\milcmarca}}}%
\parbox[t]{\dimexpr\linewidth-6.4mm\relax}{\RaggedRight #1}\par}

% Celda marcable dentro de la rubrica (tabla de autoevaluacion). Misma
% sangria francesa, pero pensada para vivir dentro de una columna X.
\newcommand{\milcopcion}[1]{%
\makebox[5.2mm][l]{\raisebox{-.55mm}{\textcolor{milcNegro}{\milcmarca}}}%
\parbox[t]{\dimexpr\linewidth-5.2mm\relax}{\RaggedRight #1}}

% ── Recursos visuales de la guía ──────────────────────────────────────
% Declarados en el bloque `recursos:` del YAML y emitidos por el builder.
% \guiaFigura   → imagen rasterizada o PDF (foto, carta celeste, montaje).
% \guiaDiagrama → fuente TikZ/circuitikz (.tex) incrustada como vector:
%                 es la vía para circuitos y esquemáticos editables.
% [#1] = texto alternativo (alt) · #2 = ruta relativa al .tex · #3 = pie
% El alt llega al PDF etiquetado (/Alt) y lo lee el lector de pantalla.
\newcommand{\guiaFigura}[3][]{%
\begin{center}
\includegraphics[width=0.88\linewidth,height=0.38\textheight,keepaspectratio,alt={#1}]{#2}\\[1.5mm]
{\footnotesize\itshape\textcolor{milcNegro!75}{#3}}
\end{center}
}

\newcommand{\guiaDiagrama}[2]{%
\begin{center}
\input{#1}\\[1.5mm]
{\footnotesize\itshape\textcolor{milcNegro!75}{#2}}
\end{center}
}

\begin{document}
\thispagestyle{empty}

% ── Portada ───────────────────────────────────────────────────────────
% Antes: un rectangulo de color a pagina completa. Se veia bien en pantalla,
% pero en la fotocopiadora del colegio se comia el toner y salia gris sucio.
% Ahora la banda ocupa el tercio superior y el resto va en blanco.
\begin{tikzpicture}[remember picture,overlay]
  \path[fill=milcGradePrimary]
    (current page.north west) rectangle ([yshift=-10.6cm]current page.north east);

  \node[anchor=north west,text=milcGradeText] at ([xshift=2.0cm,yshift=-1.55cm]current page.north west)
    {\sffamily\bfseries\fontsize{12}{14}\selectfont MOMENTO};
  \node[anchor=north west,text=milcGradeText,opacity=.85] at ([xshift=2.0cm,yshift=-2.35cm]current page.north west)
    {\sffamily\bfseries\fontsize{8.5}{10}\selectfont TERRITORIO INTERIOR · SOCIOEMOCIONAL};
  \node[anchor=north east,text=milcGradeText,opacity=.85,align=right] at ([xshift=-2.0cm,yshift=-1.55cm]current page.north east)
    {\sffamily\bfseries\fontsize{9}{12}\selectfont I.E. Sor María Juliana\\Cartago, Valle del Cauca};

  \node[anchor=west,text=milcGradeText] (gradonum) at ([xshift=1.85cm,yshift=-7.1cm]current page.north west)
    {\sffamily\bfseries\fontsize{112}{112}\selectfont 8};
  % El circulito de grado (°) solo aplica a las guias de grado («11°»); en
  % Bebras y Semillero el numero grande es un momento/modulo, no un grado.
  % Se posiciona relativo al nodo `gradonum`, no en coordenadas absolutas.
  

  \node[anchor=west,text=milcGradeText,text width=11.4cm,align=left] at ([xshift=1.5cm,yshift=0.5cm]gradonum.east)
    {
     {\sffamily\bfseries\fontsize{9}{11}\selectfont\MakeUppercase{Territorio interior · Socioemocional}}\\[3mm]
     {\sffamily\bfseries\fontsize{21}{25}\selectfont\setlength{\baselineskip}{27pt}Los grupos\\[4pt]de chat\\[4pt]sin cara y sin fecha de cierre}\\[3.5mm]
     {\sffamily\bfseries\fontsize{15}{18}\selectfont Grado 7° · Momento 8}
    };
\end{tikzpicture}

\vspace*{9.4cm}
\vfill

{\sffamily\bfseries\fontsize{8}{10}\selectfont\textcolor{milcGradeDark}{LO QUE TE LLEVAS HOY}}

\begin{tcolorbox}[enhanced,colback=white,colframe=milcNegro,arc=2mm,boxrule=1.6pt,
  left=5mm,right=5mm,top=4mm,bottom=4mm,before skip=3pt,after skip=12pt,
  fontupper=\RaggedRight\fontsize{11}{15.5}\selectfont]
Tu protocolo del grupo: las seis condiciones que te sueltan la mano aplicadas a un grupo real tuyo, la regla de las dos preguntas antes de reenviar, y tres acuerdos de convivencia redactados para proponerle al grupo.
\end{tcolorbox}

\begin{tcolorbox}[enhanced,colback=milcGris,colframe=milcGris,arc=2mm,boxrule=0pt,
  left=5mm,right=5mm,top=3.5mm,bottom=3.5mm,after skip=12pt]
{\sffamily\bfseries\fontsize{8}{10}\selectfont\textcolor{milcNegro!55}{ESTA GUÍA ES DE}}\\[3mm]
\begin{tabularx}{\linewidth}{X p{3.2cm} p{3.2cm}}
\linead{\linewidth} & \linead{3.0cm} & \linead{3.0cm}\\[-1mm]
{\fontsize{8.5}{10}\selectfont\textcolor{milcNegro!55}{Nombre completo}} &
{\fontsize{8.5}{10}\selectfont\textcolor{milcNegro!55}{Grupo}} &
{\fontsize{8.5}{10}\selectfont\textcolor{milcNegro!55}{Fecha}}\\
\end{tabularx}
\end{tcolorbox}

\vfill

\noindent\textcolor{milcLinea}{\rule{\linewidth}{1pt}}\\[2mm]
{\sffamily\bfseries\fontsize{9.5}{12}\selectfont Autor: Dr. Álvaro Cárdenas Orozco}\\[1mm]
{\sffamily\fontsize{8.5}{11}\selectfont\textcolor{milcNegro!55}{Metodología MILC · Apertura ancestral · Cuidado en línea · Triángulo Dussel-Estoicismo-Floridi}}

\newpage

\section*{Apertura: Los grupos de chat: lo que se reenvía sin cara y sin fecha de cierre}

\begin{softbox}{milcGradeDark}{milcGradeSoft}{Saber ancestral}
En \textbf{Riosucio} (Caldas), cada dos años, el pueblo se burla de sí mismo \textbf{a propósito y en público}. El Carnaval tiene una figura encargada de eso: \textbf{el matachín}, que participa desde \textbf{1915} y escribe versos y coplas donde se ríe de la corrupción de los políticos, de las escenas vergonzosas de la gente influyente y de lo que se le ocurra. Hay un tipo particular, el \textbf{decretero}: el matachín que \textbf{redacta los decretos ---textos en verso, llenos de ironía, sarcasmo y pulla, sobre los mismos vecinos--- y \textbf{los lee delante de todos}}. \textbf{Ahora fíjate en las condiciones,} porque son el asunto de esta guía. \textbf{El decretero no es anónimo:} tiene nombre y está ahí parado. \textbf{No es a escondidas:} es en la plaza, delante de aquel de quien se habla. \textbf{Y tiene fecha de cierre:} el carnaval empieza y \textbf{se acaba}, y con él el decreto. \emph{Es un pueblo entero burlándose de sí mismo con las tres condiciones que un grupo de chat no tiene: cara, plaza y fecha final.}
\end{softbox}

\begin{softbox}{milcTurquesa}{milcGris}{Saber contemporáneo}
¿Por qué la gente escribe cosas que no diría de frente? \textbf{El psicólogo John Suler} le puso nombre en 2004: el \textbf{efecto de desinhibición en línea}, y describió \textbf{seis cosas} que sueltan la mano. \textbf{Anonimato} ---no se sabe quién soy, o al menos se siente así---. \textbf{Invisibilidad} ---no me ven la cara, ni yo la de nadie---. \textbf{Asincronía} ---escribo y me voy; no me toca aguantar la reacción---. \textbf{Introyección solipsista} ---al leer, le pongo al otro la voz que yo quiero, y casi siempre le pongo la peor---. \textbf{Imaginación disociativa} ---esto es como un juego, no cuenta---. Y \textbf{minimización de la autoridad} ---aquí no hay profesor ni papá mirando--- (Suler, 2004). \textbf{Y algo importante:} Suler aclara que el efecto \textbf{tiene dos caras}. La \textbf{benigna} es real: hay gente que en línea se atreve a contar lo que le pasa, a pedir ayuda, a ser amable con un desconocido. La \textbf{tóxica} es la otra: el insulto, la amenaza, la crueldad fácil. \emph{No es que internet vuelva mala a la gente: es que quita los frenos, y sin frenos sale lo bueno y lo malo con la misma facilidad.}
\end{softbox}

\begin{softbox}{milcMostaza}{milcGris}{Pregunta-puente}
\textit{El decretero se burla \textbf{con su cara, en la plaza y con fecha de cierre}. \textbf{¿Cuál de las cosas que has reenviado este año habrías dicho igual} si te hubiera tocado leerla en voz alta, delante de todos, con esa persona presente?}
\end{softbox}

\begin{softbox}{milcVerde}{milcGris}{Saber y saber hacer}
Al terminar podrás: (1) \textbf{identificar} en un grupo real tuyo cuáles de las seis condiciones de Suler están operando; (2) \textbf{explicar} por qué una captura de pantalla cambia el dueño de un mensaje, y por qué en un grupo grande casi nadie interviene; (3) \textbf{crear} una regla propia de dos preguntas antes de reenviar; (4) \textbf{proponer} tres acuerdos de convivencia redactados para un grupo del que haces parte.
\end{softbox}



\small
\begin{tabularx}{\textwidth}{>{\bfseries}p{3.0cm}Y}
\rowcolor{milcGradePrimary!12}Autor & Dr. Álvaro Cárdenas Orozco\\
DBA & Comprende los fenómenos que afectan a su territorio, regula sus emociones ante ellos y acompaña a otros con empatía (transversal 6°--11°, articulado con la política «Escuela, territorio de vida» del MEN, Resolución 006519 de 2025, y la Ley 1620 de 2013)\\
Referentes & Primera ayuda psicológica (OMS, 2011/2012) · Directrices IASC (2007) · USGS y Servicio Geológico Colombiano · Pueblo nasa y la minga (CRIC) · Dussel · Estoicismo · Floridi\\
Producto final & Tu protocolo del grupo: las seis condiciones que te sueltan la mano aplicadas a un grupo real tuyo, la regla de las dos preguntas antes de reenviar, y tres acuerdos de convivencia redactados para proponerle al grupo.\\
\end{tabularx}
\normalsize

\puente{En el momento 4 viste la palabra que hiere. \textbf{Hoy es la misma palabra, pero con dos cambios:} nadie tiene que dar la cara y \textbf{lo dicho queda escrito}.}

\milcestacion{milcGradeDark}{Ruta MILC: así vamos a aprender}
\begin{tabularx}{\textwidth}{>{\centering\arraybackslash}X>{\centering\arraybackslash}X>{\centering\arraybackslash}X>{\centering\arraybackslash}X}
\phasecard{1}{milcVerde}{milcGris}{Escucha\\[-1mm]\small Observo, escucho y pregunto.} &
\phasecard{2}{milcTurquesa}{milcGris}{Sistematización\\[-1mm]\small Reconozco y freno.} &
\phasecard{3}{milcMagenta}{milcGris}{Praxis\\[-1mm]\small Construyo y aplico.} &
\phasecard{4}{milcVino}{milcGris}{Evaluación\\[-1mm]\small Reflexión liberadora.}
\end{tabularx}


\puente{Primero, la auditoría de un grupo tuyo de verdad: qué se manda ahí, quién manda y quién nunca dice nada.}

\milcestacion{milcVerde}{Estación 1 de 3 · Escucha}

\begin{softbox}{milcVerde}{milcGris}{Escena para escuchar}
\textbf{Actividad 1 · IDENTIFICA --- La auditoría del grupo.} \textbf{Escoge un grupo real del que hagas parte} ---el del salón, el del equipo, el de los amigos--- y \textbf{no pongas nombres: usa iniciales}. \textbf{Primera parte.} Recorre los últimos días y clasifica lo que se mandó: \emph{cosas prácticas} (tareas, horarios), \emph{chistes y memes}, \emph{cosas sobre alguien que no está en el grupo}, \emph{capturas de otras conversaciones}. ¿Cuál grupo es el más grande? \textbf{Segunda parte.} ¿\textbf{Quiénes hablan} y quiénes \textbf{nunca escriben nada}? \emph{Anota cuántos son de cada uno: la mayoría de los grupos los mueven tres o cuatro personas.} \textbf{Tercera parte, la incómoda.} ¿Se ha mandado ahí algo que \textbf{no dirías de frente} a la persona de la que se trataba? ¿Lo reenviaste, te reíste o te quedaste callado? \emph{Y una última: ¿alguna vez te enteraste de que algo que tú escribiste terminó en otro grupo?}
\end{softbox}

{\sffamily\bfseries\fontsize{8}{10}\selectfont\textcolor{milcNegro!55}{MARCA CUANDO LO HAYAS HECHO}}
\milccheckrow{Clasificaste lo que se manda en el grupo en las cuatro categorías y sabes cuál pesa más.}
\milccheckrow{Contaste cuántos escriben y cuántos nunca escriben.}
\milccheckrow{Anotaste tu propia participación ---reenviar, reírte o callar--- sin justificarte.}

\begin{infoband}{milcVerde}


\textbf{En tu cuaderno (Actividad 1 · IDENTIFICA).} Título: «Actividad 1. Auditoría del grupo». \textbf{Formato:} las cuatro categorías con cuántos mensajes de cada una + quiénes hablan y quiénes no + tu parte. \textbf{Extensión:} media página. \textbf{Sabes que terminaste cuando} puedes decir \textbf{qué porcentaje del grupo lo mueven tres personas}. Esta página es tuya: \textbf{nadie más tiene que leerla si tú no quieres}.
\end{infoband}

\puente{Ya lo tienes. Ahora, las seis cosas que sueltan la mano, qué hace una captura y por qué entre más gente hay, menos gente interviene.}

\milcestacion{milcTurquesa}{Fase 2 · Sistematización: entender el grupo}

\begin{softbox}{milcTurquesa}{milcGris}{Cuatro claves sobre lo que pasa en un grupo}
Cuatro claves para entender por qué en un grupo pasa lo que no pasaría en el salón, con las mismas personas.
\end{softbox}

\milcpaso{1}{milcTurquesa}{\textbf{Las seis cosas que sueltan la mano.} Suler las nombró: \textbf{anonimato}, \textbf{invisibilidad}, \textbf{asincronía}, \textbf{introyección solipsista} ---leer poniéndole al otro la voz que uno quiere---, \textbf{imaginación disociativa} ---«esto no cuenta, es solo el chat»--- y \textbf{minimización de la autoridad} (Suler, 2004). \emph{Míralas juntas y se entiende todo:} el que escribe no ve la cara del otro, no le toca aguantar la reacción, se imagina el tono que le conviene y siente que nada de eso es del todo real. \textbf{Y la más engañosa es la cuarta:} cuando lees un mensaje sin signos, \textbf{tú} le pones el tono ---y si estás de mal genio, le pones el peor---. \emph{La mitad de las peleas de grupo empiezan ahí, no en lo que el otro quiso decir.}}
\milcpaso{2}{milcTurquesa}{\textbf{La captura de pantalla cambia el dueño del mensaje.} Esto es lo específico de hoy y no tiene equivalente en el recreo: \textbf{cualquier cosa que escribas puede ser sacada de donde la dijiste y puesta en otra parte}, sin el antes ni el después, sin el tono, sin quién había preguntado qué. \emph{Una frase que en su conversación se entendía, sola parece otra cosa ---y a veces parece exactamente lo contrario---.} \textbf{Dos consecuencias prácticas.} \emph{Primera:} lo que escribes ya no es tuyo desde el momento en que lo mandas. \emph{Segunda, y esta es la que casi nadie piensa:} \textbf{la captura de una conversación privada expone a los dos}, también al que no hizo nada.}
\milcpaso{3}{milcTurquesa}{\textbf{Cuantos más son, menos interviene alguien.} En un grupo de cuarenta, cuando alguien manda algo cruel, \textbf{casi nadie dice nada} ---y no es porque a los cuarenta les parezca bien---. \emph{Cada uno piensa lo mismo:} «alguien más va a decir algo», «no es conmigo», «si hablo me caen a mí». \textbf{Y como todos piensan igual, no habla ninguno, y el silencio termina pareciendo aprobación.} \emph{Aquí vale lo mismo que en el momento 4:} el público es el que hace funcionar la cosa. \textbf{Pero en el chat hay una salida que el salón no tiene:} \emph{no hace falta responderle al grupo}. Un mensaje privado a la persona afectada ---«vi lo que mandaron, qué mala jugada»--- cambia el día de alguien \textbf{y no le cuesta nada al que lo manda}.}
\milcpaso{4}{milcTurquesa}{\textbf{Lo que el carnaval tiene y el grupo no: fecha de cierre.} El decretero de Riosucio se burla \textbf{con nombre, en la plaza y en unos días que se acaban}. \emph{Un grupo de chat no tiene ninguna de las tres:} no hay que dar la cara, no hay plaza y \textbf{no hay último día} ---el mensaje de hace dos años sigue ahí, buscable---. \textbf{Por eso funciona la prueba de la plaza:} antes de mandar algo, pregúntate si lo leerías en voz alta, con tu nombre, delante de esa persona. \emph{Ojo, no es que la burla esté prohibida:} Riosucio lleva un siglo burlándose de sus propios políticos. \textbf{Es que allá la burla tiene autor, tiene público y tiene final.}}

\begin{drawbox}{milcTurquesa}{Las dos preguntas antes de reenviar}
\begin{pasoslista}
\item \textbf{¿Yo pondría mi nombre en esto?} \emph{Es la prueba de la plaza: si la gracia depende de que no se sepa quién fue, ya sabes qué es.}
\item \textbf{¿Esto es mío para mandarlo?} Una captura de una conversación privada \textbf{no lo es} ---y expone también al otro---.
\item \textbf{Si alguna respuesta es «no», no lo mandes.} No hace falta explicar nada: no mandarlo no se nota.
\item \textbf{Y si ya lo mandaste,} el momento 6 tiene el procedimiento: reconocer, ofrecer arreglo, corregir \textbf{donde} se hizo el daño.
\end{pasoslista}
\end{drawbox}

\begin{softbox}{milcMostaza}{milcGris}{Errores comunes y su corrección}
\textbf{«Yo solo lo reenvié»:} reenviar es publicar; el que reenvía es el que lo lleva a gente nueva. \textbf{«Es que era privado»:} no existe lo privado con captura. \textbf{«Yo no dije nada»:} en un grupo, el silencio de muchos parece aprobación. \textbf{Contestarle al grupo:} casi siempre alarga el espectáculo; el privado sirve más. \textbf{Salirse y ya:} te protege a ti y deja a la persona sola. \textbf{Leer con la peor voz posible:} la mitad de las peleas empiezan en un tono que le pusiste tú al mensaje. \textbf{Guardar capturas «por si acaso»:} salvo que sea para reportar un caso serio a un adulto, es munición.
\end{softbox}

\begin{infoband}{milcTurquesa}


\textbf{En tu cuaderno (Actividad 2 · EXPLICA).} Título: «Actividad 2. Las seis condiciones». \textbf{Formato:} las seis de Suler, y al frente de cada una si opera o no en el grupo que auditaste, con un ejemplo concreto de esa condición. \textbf{Extensión:} media página. \textbf{Sabes que terminaste cuando} puedes explicar por qué \textbf{una captura de una conversación privada expone también al que no hizo nada}.
\end{infoband}

\puente{Ya sabes qué pasa. Es hora de ponerle reglas al grupo ---empezando por una tuya, de dos preguntas, que se aplica en tres segundos---.}

\milcestacion{milcMagenta}{Fase 3 · Praxis: poner reglas al grupo}

\begin{softbox}{milcMagenta}{milcGris}{Cómo se le ponen reglas a un grupo que no las tiene}
Ahora se convierte en reglas. Primero una regla tuya, que se aplica en tres segundos y no hay que negociar con nadie; después, tres acuerdos para proponerle al grupo.
\end{softbox}

\milcpaso{1}{milcMagenta}{\textbf{La auditoría, terminada (8 min).} Vuelve a tu grupo y marca \textbf{cuáles de las seis condiciones de Suler} están operando ahí, con un ejemplo de cada una. \emph{Vas a encontrar que la asincronía y la invisibilidad están siempre, y que el anonimato casi nunca ---en un grupo todos saben quién eres---.} \textbf{Eso es interesante:} quiere decir que lo que suelta la mano \textbf{no es el anonimato}, sino no tener que aguantar la cara del otro.}
\milcpaso{2}{milcMagenta}{\textbf{Las dos preguntas (7 min).} Escribe tu regla en una tarjeta o en una nota del celular: \textbf{«¿pondría mi nombre?»} y \textbf{«¿es mío para mandarlo?»}. \emph{Que sea de dos preguntas y no de diez tiene una razón:} en el momento real uno tiene tres segundos y el dedo encima. \textbf{Pruébala con cinco cosas} que hayas reenviado esta semana ---incluidas las inofensivas---.}
\milcpaso{3}{milcMagenta}{\textbf{El mensaje privado (10 min, en parejas).} Escribe el mensaje que le mandarías \textbf{en privado} a alguien de quien se estuvieran burlando en un grupo. \textbf{Máximo dos frases}, sin sermón y sin pedirle que haga nada. \emph{Compáralo con el de tu pareja de trabajo:} van a ver que los que sirven se parecen ---dicen que lo vieron, dicen que estuvo mal y no piden nada a cambio--- y que los que no sirven traen consejo adentro.}
\milcpaso{4}{milcMagenta}{\textbf{Los tres acuerdos (10 min, en grupos).} Redacten \textbf{tres acuerdos} para un grupo real, con dos condiciones: que se puedan \textbf{comprobar} y que \textbf{no sean solo prohibiciones}. \emph{Ejemplos de acuerdos que sí funcionan:} «lo de otra conversación no se trae aquí» · «si alguien pide que borren algo suyo, se borra sin discutir» · «el que se sienta aludido lo dice y se para ahí». \textbf{Y el que más cuesta y más sirve:} \emph{«el que quiera pelear, lo habla por fuera»}. \textbf{Escriban también qué pasa cuando alguien lo rompe:} un acuerdo sin eso dura una semana.}

\begin{drawbox}{milcMagenta}{Antes de cerrar tu protocolo}
\begin{checklista}
\item Marcaste cuáles de las \textbf{seis condiciones} operan en tu grupo, con un ejemplo de cada una.
\item Tu regla de las dos preguntas está \textbf{escrita donde la vas a ver} ---tarjeta o nota del celular---.
\item La probaste con cinco cosas que reenviaste esta semana.
\item Tienes escrito un \textbf{mensaje privado} de dos frases, sin sermón y sin pedir nada.
\item Los tres acuerdos se pueden \textbf{comprobar} y no son solo prohibiciones.
\item Escribieron \textbf{qué pasa cuando alguien rompe un acuerdo}.
\end{checklista}
\end{drawbox}

\begin{softbox}{milcMagenta}{milcGris}{Plantilla de trabajo}
\textbf{Mi regla, en dos preguntas.} \emph{«¿Pondría mi nombre en esto? · ¿Es mío para mandarlo?»} \\[1mm]
\textbf{El mensaje privado.} \emph{«Vi lo que mandaron en el grupo. Me pareció una mala jugada.»} ---dos frases, sin consejo---. \\[1mm]
\textbf{Los tres acuerdos.} \emph{1) \_\_\_\_ · 2) \_\_\_\_ · 3) \_\_\_\_}, y \emph{«si alguien lo rompe, \_\_\_\_»}.
\end{softbox}

\begin{infoband}{milcMagenta}


\textbf{En tu cuaderno (Actividad 3 · CREA).} Título: «Actividad 3. Mi protocolo del grupo». \textbf{Formato:} las seis condiciones marcadas + tu regla de dos preguntas probada con cinco casos + el mensaje privado + los tres acuerdos con su consecuencia. \textbf{Extensión:} una página. \textbf{Sabes que terminaste cuando} tus tres acuerdos \textbf{se pueden comprobar} ---es decir, cuando cualquiera del grupo podría decir si se cumplieron o no---.
\end{infoband}

\puente{El producto es un protocolo. \emph{No es una lista de prohibiciones: es lo que uno decide antes, para no tener que decidirlo con el dedo encima del botón.}}

\milcestacion{milcMostaza}{Producto de la sesión}
\textbf{Protocolo del grupo: las seis condiciones aplicadas, la regla de las dos preguntas y tres acuerdos comprobables}.
\begin{softbox}{milcMostaza}{milcGris}{Criterios de calidad}
(1) Las seis condiciones de Suler están identificadas en un grupo real, con ejemplo concreto de cada una. (2) La regla de las dos preguntas está escrita y probada con casos propios de la semana. (3) El mensaje privado cabe en dos frases, no da consejos y no pide nada a cambio. (4) Los tres acuerdos son comprobables y no se limitan a prohibir. (5) El protocolo dice qué se hace cuando alguien rompe un acuerdo, no solo qué no se debe hacer.
\end{softbox}

\puente{Antes de cerrar, revisa algo: si tus acuerdos solo dicen «no hacer», les falta la mitad. ¿Qué se hace cuando alguien ya lo hizo?}

\milcestacion{milcVino}{Evaluación liberadora · cinco dimensiones}

\begin{softbox}{milcVino}{milcGris}{Sentido de la evaluación}
Evaluar no es solo revisar una nota. Es reconocer qué aprendí, qué cambió en mí, qué emociones manejé, cómo me relaciono con la ciudad y la comunidad, y qué le dejaré a quien viene detrás.
\end{softbox}

\textbf{Mi reflexión liberadora en cinco dimensiones.} Completa cada frase iniciada:

\small
\begin{tabularx}{\textwidth}{>{\bfseries\RaggedRight\arraybackslash}p{3.4cm}Y>{\RaggedRight\arraybackslash}p{4.6cm}}
\rowcolor{milcVino}\color{white}Dimensión & \color{white}\bfseries Mi respuesta (frame starter) & \color{white}\bfseries Algo que descubrí\\
1. Personal & Lo que más cambió en mí fue \linead{6.5cm} & \linead{4.2cm}\\
2. Emocional & La emoción más fuerte fue \linead{2.5cm} y la regulé \linead{3cm} & \linead{4.2cm}\\
3. Ciudadana & En democracia esto importa porque \linead{6cm} & \linead{4.2cm}\\
4. Local (Cartago) & En mi barrio sería útil para \linead{6.5cm} & \linead{4.2cm}\\
5. Intergeneracional & A mi abuela/abuelo le diría \linead{6.5cm} & \linead{4.2cm}\\
\end{tabularx}
\normalsize

\begin{infoband}{milcVino}
\textbf{Para la próxima sustentación.} Vuelve a esta tabla en seis meses. Verás que las respuestas cambian — no porque lo aprendido cambie, sino porque tú cambias. La evaluación liberadora es una conversación contigo a través del tiempo.
\end{infoband}


\puente{Cerramos con tres voces: la cara que no está del otro lado, la demora como remedio, y qué se cultiva o se arruina en un espacio compartido.}

\milcestacion{milcGradeDark}{Triángulo de pensamiento · cierre filosófico}

\begin{softbox}{milcVino}{milcGris}{Dussel · lente del nosotros}
\textit{«El otro se revela realmente como otro\ldots como el pobre, el oprimido; el que, a la vera del camino, fuera del sistema, muestra su rostro sufriente y sin embargo desafiante: «¡Tengo hambre!, ¡tengo derecho a comer!»»}

{\footnotesize\itshape\color{milcNegro!70}--- Enrique Dussel · Filosofía de la liberación (1977)}

Dussel insiste en una palabra: \textbf{rostro}. \emph{Y el rostro es exactamente lo que un grupo de chat borra.} \textbf{Del otro lado no hay una cara: hay un nombre en una lista}, y a un nombre en una lista se le puede escribir cualquier cosa. \textbf{Fíjate en lo que dice Suler sobre la invisibilidad y la asincronía:} el que escribe \textbf{no ve la reacción} y \textbf{no tiene que quedarse a aguantarla}. \emph{Es la misma idea, dicha con setenta años de diferencia:} donde no hay rostro, el freno no aparece. \textbf{Y hay una consecuencia práctica:} si vas a escribir algo pesado, \textbf{ponle la cara} ---dilo de frente o no lo digas---.

\textbf{¿Le he escrito a alguien algo que no le habría dicho teniéndolo enfrente?}
\end{softbox}
\linead{\linewidth}\\[1mm]
\linead{\linewidth}

\begin{softbox}{milcMostaza}{milcGris}{Estoicismo · Séneca · Sobre la ira, libro II (c. 45 d.C.) · cuidado interior}
\textit{«El mayor remedio para la ira es la demora.»}

{\footnotesize\itshape\color{milcNegro!70}--- Séneca · Sobre la ira, libro II (c. 45 d.C.)}

Séneca no dice que la ira esté mal: dice que \textbf{el remedio es esperar}. \emph{Y el chat está diseñado exactamente para lo contrario:} el mensaje sale en un segundo, y el segundo siguiente ya no se puede recoger. \textbf{La versión práctica de esta frase es ridículamente simple y funciona:} cuando algo te dé rabia en un grupo, \textbf{escribe la respuesta y no la mandes}. Déjala ahí. \emph{Vuelve en diez minutos y vas a borrar el ochenta por ciento} ---y lo que quede, que casi siempre son dos líneas, sí valía la pena---. \textbf{La demora no es tragarse las cosas:} es la diferencia entre responder y reaccionar.

\textbf{¿Cuál fue el último mensaje que mandé con rabia, y qué habría escrito diez minutos después?}
\end{softbox}
\linead{\linewidth}\\[1mm]
\linead{\linewidth}

\begin{softbox}{milcTurquesa}{milcGris}{Floridi · lente de la infoesfera}
\textit{«El florecimiento de las entidades informacionales y de toda la infoesfera debe promoverse preservando, cultivando y enriqueciendo sus propiedades.»}

{\footnotesize\itshape\color{milcNegro!70}--- Luciano Floridi · La ética de la información (2013; trad.)}

Un grupo de chat es \textbf{un lugar compartido}, aunque no lo parezca: los cuarenta viven ahí todos los días. \emph{Y a un lugar compartido le pasa lo mismo que a un salón o a una cuadra:} se cuida o se arruina, y \textbf{casi nunca se arruina de golpe} ---se arruina un mensaje a la vez, hasta que a los callados les da pereza abrirlo---. \textbf{Por eso los acuerdos importan más que las prohibiciones:} prohibir es quitar; \emph{cultivar} es decidir qué sí se hace ahí. \textbf{Y esto conecta con el momento 6:} cuando el daño fue público, la reparación también tiene que serlo ---corregir en el mismo grupo donde se dijo---.

\textbf{De los grupos en los que estoy, ¿en cuál me da pereza entrar? ¿Desde cuándo?}
\end{softbox}
\linead{\linewidth}\\[1mm]
\linead{\linewidth}

\begin{infoband}{milcNegro}
\textbf{Nota para el docente.} Las tres son citas textuales y llevan su referencia debajo. Puedes remitir a la obra en clase.
\end{infoband}

\milcestacion{milcVino}{Mi autoevaluación · marca una casilla por fila}

{\small
\setlength{\extrarowheight}{2.2pt}
\begin{tabularx}{\textwidth}{>{\RaggedRight\arraybackslash\bfseries}p{3.0cm}YYY}
\rowcolor{milcVino}\color{white}\bfseries Criterio & \color{white}\bfseries Lo logré & \color{white}\bfseries En proceso & \color{white}\bfseries Necesito apoyo\\[.6mm]
Comprensión del tema & \milcopcion{Explico las ideas con mis palabras.} & \milcopcion{Reconozco las ideas, falta precisión.} & \milcopcion{Necesito repasar lo básico.}\\[.8mm]
\rowcolor{milcGris}Calidad del protocolo & \milcopcion{Tengo tres acuerdos concretos y sé qué hacer antes de reenviar algo.} & \milcopcion{Reconozco cuándo un grupo se pasa, pero solo se me ocurre salirme.} & \milcopcion{Sigo pensando que lo que pasa en el grupo se queda en el grupo.}\\[.8mm]
Aplicación y producto & \milcopcion{Mi producto cumple los criterios.} & \milcopcion{Está armado, con ajustes pendientes.} & \milcopcion{Aún está en construcción.}\\[.8mm]
\rowcolor{milcGris}Reflexión 5 dimensiones & \milcopcion{Respondí cada una con honestidad.} & \milcopcion{Respondí algunas con detalle.} & \milcopcion{Mis respuestas son muy generales.}\\[.8mm]
Triángulo de pensamiento & \milcopcion{Conecté las 3 voces con mi trabajo.} & \milcopcion{Conecté al menos una.} & \milcopcion{No las relaciono con mi proceso.}\\[.8mm]
\end{tabularx}}

\vspace{3mm}
\textbf{Hoy entendí que:}\\[1mm]
\linead{\linewidth}\\[1mm]
\linead{\linewidth}

\textbf{Mi compromiso para la próxima sesión es:}\\[1mm]
\linead{\linewidth}\\[1mm]
\linead{\linewidth}

\begin{softbox}{milcMostaza}{milcGris}{Cita de despedida · Marco Aurelio}
\textit{``El obstáculo en el camino se vuelve el camino. Lo que estorba al camino se vuelve camino.''}\\[1mm]
Cada error de hoy, bien observado, se convierte en pieza del camino que pisarás mañana.
\end{softbox}

\small
\begin{tabularx}{\textwidth}{>{\bfseries}p{3.6cm}Y}
\rowcolor{milcGradeDark!12}Nombre completo: & \linead{10cm}\\
Fecha: & \linead{4cm} \quad Curso/grupo: \linead{4cm}\\
Mi firma: & \linead{9cm}\\
\end{tabularx}
\normalsize

\begin{infoband}{milcGradeDark}
\textbf{Para la próxima sesión.} Dos cosas. \textbf{Primero, aplica la regla una vez} ---una sola: algo que ibas a reenviar y no reenviaste---, y anota qué era y qué pasó por no mandarlo. \emph{Vas a descubrir el hallazgo más útil de esta guía: no pasó nada. Nadie lo notó.} \textbf{Y si te alcanza el valor, propón un acuerdo} en un grupo tuyo, aunque sea el más chiquito. \textbf{Segundo, pregunta en tu casa por la burla que sí se firmaba:} averigua con alguien mayor \textbf{si en su pueblo o en su barrio había coplas, versos, apodos de carnaval o «decretos»} ---quién los hacía, si se decían delante de la persona, si alguien se enojaba y qué pasaba después---. \textbf{Trae:} lo que no reenviaste y lo que te contaron. \emph{Vas a encontrar que la burla es viejísima y que casi siempre tenía tres cosas que hoy no tiene: autor conocido, público presente y una fecha en que se acababa.}
\end{infoband}

\end{document}
